\documentclass[10pt,a4paper]{amsart}
\usepackage[margin=19mm]{geometry}
\usepackage{amsmath,amssymb,mathtools}
\usepackage[scr=rsfs,cal=euler]{mathalfa}
\usepackage{xfrac}
\usepackage[unicode,hidelinks,bookmarks=false]{hyperref}
\newcommand{\ie}{i.e.\ }
\newcommand{\cf}{cf.\ }
\newcommand{\resp}{respectively\ }
\newcommand{\Subsection}{\subsection}
\providecommand{\nobreakdash}{\nobreak\hbox{-}}
\setlength{\emergencystretch}{2em}
\multlinegap0pt
\usepackage{bm}
\usepackage{bbm}
\usepackage{dsfont}
\usepackage{xspace}
\usepackage{cancel}
\usepackage{mathtools}
\usepackage{amsthm}
\makeatletter
\@ifundefined{theorem}{\newtheorem{theorem}{Theorem}}{}
\makeatother
\title[On unbiased estimators for functions of the rate parameter o]{On unbiased estimators for functions of the rate parameter of the exponential distribution}
\author{Roberto Vila, Eduardo Yoshio Nakano}
\date{Source: arXiv:2506.20005v2 (CC BY 4.0)}
\begin{document}
\maketitle

\noindent\emph{Source and licence.} On unbiased estimators for functions of the rate parameter of the exponential distribution. Version arXiv:2506.20005v2 (CC BY 4.0), distributed under the Creative Commons Attribution 4.0 International licence, as recorded in the arXiv licence metadata for this exact version and shown on the abstract page https://arxiv.org/abs/2506.20005v2. Full rights notice: licenses/arxiv-2506.20005-CC-BY-4.0.txt.\par\smallskip

\noindent\textit{Editorial scope.} Counted unit: the Theorem whose source label is \texttt{main theorem} in the article (source0.tex lines 275--287, source label \texttt{main theorem}), delivered together with the article's own complete original proof of it (source0.tex lines 289--397, one \texttt{proof} environment of 109 lines). No mathematics is written, repaired or re-derived: statement and proof are the article's own text line for line. Required context retained uncounted: eq-1 (lines 258--262). Every definition, display and label of those lines is retained unchanged. Omitted from this package and not redistributed: 1--257, 263--274, 288--288, 398--1537. Nothing inside a retained excerpt is omitted or altered: every retained body is delivered exactly as the article's lines are written, so no omission receipt of any kind accompanies this package. Citations: the retained excerpts carry no citation command at all, so no bibliography entry of the article is transcribed and no reference is redistributed. Reference adaptation: none was needed and none was made; every reference inside the delivered unit resolves inside it, so no unresolved reference or citation is produced. Printed slips kept literal: no printed word, symbol, hypothesis or number of the counted statement or of its proof was altered. Layout and class: the article's own class is \texttt{article} with packages including amsmath, graphicx, psfrag, epsf, booktabs, enumerate, url, amssymb, amsfonts, amsthm, mathtools, mathrsfs, xcolor, bm; the wrapper is the lane's pinned \texttt{amsart} editorial wrapper at 10pt on A4, which restates the theorem-like environments the retained text needs and adds the article's own macro definitions that the retained text needs (none), with extra wrapper packages (bm, bbm, dsfont, xspace, cancel, mathtools). The delivered numbering is the unit's own. Assets: no retained span contains a figure environment, a graphics-inclusion command, a PDF-page inclusion, a table or a TikZ picture; no image file is copied into this package, the package declares no asset dependency and no image file is redistributed with it. Licence: version arXiv:2506.20005v2 of this article is distributed under the Creative Commons Attribution 4.0 International licence, as recorded in the arXiv licence metadata for this exact version and shown on the abstract page https://arxiv.org/abs/2506.20005v2. A full rights notice is delivered at licenses/arxiv-2506.20005-CC-BY-4.0.txt. Characters: every retained excerpt is pure ASCII, so the delivered engine, XeLaTeX with its default Latin Modern text font, reports no missing character.
\begin{align}\label{eq-1}
	\mathbb{E}\left[\phi(\overline{X})\right]
	=
	\xi(\lambda).
\end{align}

\begin{theorem}\label{main theorem}
	If $X_1,\ldots,X_n$ is a random sample of size $n$ from $X\sim\exp(\lambda)$, then, the following sample function
	\begin{align}\label{final-exp}
		\phi(\overline{X})
		=
		\int_0^{\infty} 
		\mathds{1}_{\{\overline{X}\geqslant v\}}
		\left(1-{v\over \overline{X}}\right)^{n-1} 
		\mathscr{L}^{-1}\left\{\xi\left({s\over n}\right)
		\right\}(v){\rm d}v
	\end{align}
	is an unbiased estimator for $\xi(\lambda)$, where $F(s)=\mathscr{L}\left\{f(x)\right\}(s)=\int_0^\infty f(x)\exp(-sx){\rm d}x$ denotes the Laplace transform and $f(x)=\mathscr{L}^{-1}\left\{F(s)\right\}(x)$ is its respective inverse.  In the above we are assuming that the inverse of the Laplace transform and the respective improper integral exist.
\end{theorem}

\begin{proof}
	It is well-known that  $\overline{X}\sim\text{Gamma}(n,n\lambda)$. Then the equality \eqref{eq-1} can be written as
	\begin{align*}
		{(n\lambda)^n\over \Gamma(n)}
		\int_{0}^{\infty}
		\phi(x)
		x^{n-1}
		\exp(-n\lambda x)
		{\rm d}x
		=
		\xi(\lambda),
	\end{align*}
	which can be expressed in function of Laplace transform as follows
	\begin{align*}
		\mathscr{L}\left\{	\phi(x)
		x^{n-1}\right\}(n\lambda)
		=
		\Gamma(n)\,
		{\displaystyle \xi\left({1\over n} \, (n\lambda)\right)\over (n\lambda)^n}
		=
		\Gamma(n)\,
		{\displaystyle \xi\left({s\over n}\right)\over s^n}
		\Bigg\vert_{s=n\lambda}.
	\end{align*}
	Applying the inverse Laplace transform (if it exists) to both sides gives
	\begin{align}\label{main-identit-0}
		\phi(x)
		x^{n-1}
		=
		\Gamma(n)\,
		\mathscr{L}^{-1}\left\{	
		{\displaystyle \xi\left({s\over n}\right)\over s^n}
		\right\}
		(x)
		=
		\Gamma(n)\,
		\mathscr{L}^{-1}\left\{	
		F(s)H(s)
		\right\}
		(x),
	\end{align}	
	%Thus, an unbiased estimator for $\theta(\lambda)$ is given by
	%\begin{align}\label{main-identity}
	%	g(\overline{X})
	%	=
	%	\Gamma(n)\,
	%	{1\over \overline{X}^{n-1}} \, 
	%	\mathscr{L}^{-1}\left\{	
	%	{\displaystyle \theta\left({s\over n}\right)\over s^n}
	%	\right\}
	%	(\overline{X}),
	%\end{align}
	%where $\overline{X}=(1/n)\sum_{i=1}^{n}X_i$ is the sample mean.
	where we are adopting the notation $F(s)=1/s^n$ and $H(s)=\xi\left({s/n}\right)$. 
	
	Since
	\begin{align*}
		f(x)=  \mathscr{L}^{-1}\left\{	F(s)
		\right\}(x)
		=
		{x^{n-1}\over\Gamma(n)}
		\quad
		\text{and}
		\quad 
		h(x)=\mathscr{L}^{-1}\left\{H(s)
		\right\}(x)
		=
		\mathscr{L}^{-1}\left\{\xi\left({s\over n}\right)
		\right\}(x),
	\end{align*}
	by applying the convolution theorem for the Laplace transform, we obtain
	\begin{align}
		\mathscr{L}^{-1}\left\{	
		F(s) H(s)
		\right\}
		(x)
		=
		(f*h)(x)
		&=
		\int_0^x f(x-v)h(v){\rm d}v
		\nonumber
		\\[0,2cm]
		&=
		{1\over\Gamma(n)}
		\int_0^x (x-v)^{n-1} 
		\mathscr{L}^{-1}\left\{\xi\left({s\over n}\right)
		\right\}(v){\rm d}v
		\label{ide-3-3-1}
		\\[0,2cm]
		&=
		{x^{n-1} \over\Gamma(n)}
		\int_0^\infty 
		\mathds{1}_{\{x\geqslant v\}}
		\left(1-{v\over x}\right)^{n-1} 
		\mathscr{L}^{-1}\left\{\xi\left({s\over n}\right)
		\right\}(v){\rm d}v, \label{ide-3-3}
	\end{align}
	where $x\mapsto(f*h)(x)$ denotes  the convolution function. By plugging \eqref{ide-3-3} in \eqref{main-identit-0}, we get
	\begin{align}
		\phi(x)
		=
		\int_0^\infty 
		\mathds{1}_{\{x\geqslant v\}}
		\left(1-{v\over x}\right)^{n-1} 
		\mathscr{L}^{-1}\left\{\xi\left({s\over n}\right)
		\right\}(v){\rm d}v.
	\end{align}
	Thus, the unbiased estimator $\phi(\overline{X})$ takes the form given in \eqref{final-exp}, which completes the proof of the theorem.
\end{proof}

\end{document}
